\documentclass[11pt,a4paper]{article}

% ---- Packages ----
\usepackage[utf8]{inputenc}
\usepackage[T1]{fontenc}
\usepackage{amsmath,amssymb,amsthm}
\usepackage{mathtools}
\usepackage[margin=1in,headheight=14pt]{geometry}
\usepackage{hyperref}
\usepackage{enumitem}
\usepackage{xcolor}
\usepackage{tcolorbox}
\usepackage{fancyhdr}
\usepackage{booktabs}

% ---- Hyperref ----
\hypersetup{
    colorlinks=true,
    linkcolor=red,
    citecolor=blue,
    urlcolor=blue
}

% ---- Header ----
\pagestyle{fancy}
\fancyhf{}
\fancyhead[L]{\textit{Linear Algebra Review}}
\fancyhead[R]{\thepage}
\renewcommand{\headrulewidth}{0.4pt}

% ---- Theorem Environments ----
\theoremstyle{definition}
\newtheorem{definition}{Definition}[section]
\newtheorem{example}{Example}[section]
\newtheorem{exercise}{Exercise}[section]

\theoremstyle{plain}
\newtheorem{theorem}{Theorem}[section]
\newtheorem{lemma}[theorem]{Lemma}
\newtheorem{proposition}[theorem]{Proposition}
\newtheorem{corollary}[theorem]{Corollary}

\theoremstyle{remark}
\newtheorem{remark}{Remark}[section]

% ---- Colored Boxes ----
\tcbuselibrary{theorems,skins,breakable}

\newtcolorbox{keyresult}[1][]{
    colback=green!5!white,
    colframe=green!50!black,
    fonttitle=\bfseries,
    title=#1,
    breakable
}

\newtcolorbox{rlconnection}[1][]{
    colback=blue!5!white,
    colframe=blue!50!black,
    fonttitle=\bfseries,
    title=#1,
    breakable
}

\newtcolorbox{warning}[1][]{
    colback=red!5!white,
    colframe=red!50!black,
    fonttitle=\bfseries,
    title=#1,
    breakable
}

% ---- Operators ----
\newcommand{\R}{\mathbb{R}}
\newcommand{\C}{\mathbb{C}}
\newcommand{\F}{\mathbb{F}}
\newcommand{\N}{\mathbb{N}}
\DeclareMathOperator{\rank}{rank}
\DeclareMathOperator{\nullity}{nullity}
\DeclareMathOperator{\range}{range}
\DeclareMathOperator{\nullspace}{null}
\DeclareMathOperator{\tr}{tr}
\DeclareMathOperator{\diag}{diag}
\DeclareMathOperator{\Span}{span}

% ---- Title ----
\title{\textbf{Linear Algebra Review}\\[10pt]
\large Vector Spaces, Linear Maps, and the Spectral Theorem\\[5pt]
\normalsize Phase 0: Prerequisites Review $\mid$ Document 2 of 3\\
\normalsize Primary texts: Axler, \emph{Linear Algebra Done Right}; Hoffman \& Kunze, \emph{Linear Algebra}}
\author{Curriculum: A Rigorous Learning Path to Multi-Agent Deep RL}
\date{February 2026}

\begin{document}

\maketitle

\begin{abstract}
This document provides a rigorous refresher of finite-dimensional linear algebra,
focused on the results needed for functional analysis, optimization theory, and
machine learning. We develop vector spaces, linear maps, determinants, eigenvalue
theory, inner product spaces, the spectral theorem for real symmetric matrices, the
singular value decomposition, and matrix norms. Every theorem is stated precisely
and proved in full. The reader who completes this review will have the algebraic
foundations required for Phase~01 (functional analysis, optimization) and beyond.
\end{abstract}

\tableofcontents
\newpage

%% ============================================================
%% SECTION 1: INTRODUCTION
%% ============================================================
\section{Introduction}
\label{sec:intro}

This document is a refresher, not a first course. We assume the reader has seen
linear algebra before but needs to recall the key definitions, theorems, and proof
techniques. We are selective: only results that the rest of the curriculum depends
on are included.

\medskip\noindent\textbf{What this document covers.}
\begin{enumerate}[nosep]
    \item \textbf{Vector spaces} --- axioms, subspaces, span, linear independence, basis, dimension.
    \item \textbf{Linear maps} --- null space, range, rank-nullity theorem, matrix representation.
    \item \textbf{Determinants} --- properties, multiplicativity, invertibility criterion.
    \item \textbf{Eigenvalues and eigenvectors} --- characteristic polynomial, diagonalizability.
    \item \textbf{Inner product spaces} --- Cauchy--Schwarz, Gram--Schmidt, orthogonal projection.
    \item \textbf{The spectral theorem} --- real symmetric matrices are orthogonally diagonalizable.
    \item \textbf{Singular value decomposition} --- existence, geometric interpretation, low-rank approximation.
    \item \textbf{Matrix norms and the matrix exponential} --- operator norm, Frobenius norm, $e^A$.
\end{enumerate}

\begin{rlconnection}[Why Linear Algebra Matters for AI/ML]
Linear algebra is the computational backbone of machine learning. Functional analysis
(Phase~01) generalizes vector spaces to infinite dimensions. Optimization theory uses
eigenvalues and positive definiteness to characterize convex functions. Neural networks
are compositions of linear maps with pointwise nonlinearities. The spectral theorem
underlies PCA, covariance analysis, and the Fisher information matrix in natural policy
gradient methods.
\end{rlconnection}

%% ============================================================
%% SECTION 2: VECTOR SPACES
%% ============================================================
\section{Vector Spaces}
\label{sec:vector_spaces}

Throughout, $\F$ denotes either $\R$ or $\C$.

\begin{definition}[Vector Space]
\label{def:vector_space}
A \emph{vector space} over a field $\F$ is a set $V$ together with operations
of addition $+: V \times V \to V$ and scalar multiplication $\cdot: \F \times V \to V$
satisfying:
\begin{enumerate}[nosep]
    \item \textbf{Commutativity:} $u + v = v + u$ for all $u, v \in V$.
    \item \textbf{Associativity of addition:} $(u + v) + w = u + (v + w)$ for all $u, v, w \in V$.
    \item \textbf{Additive identity:} There exists $0 \in V$ such that $v + 0 = v$ for all $v \in V$.
    \item \textbf{Additive inverse:} For each $v \in V$, there exists $w \in V$ with $v + w = 0$.
    \item \textbf{Multiplicative identity:} $1 \cdot v = v$ for all $v \in V$.
    \item \textbf{Associativity of scalar multiplication:} $a(bv) = (ab)v$ for all $a, b \in \F$, $v \in V$.
    \item \textbf{Distributive laws:} $a(u + v) = au + av$ and $(a + b)v = av + bv$.
\end{enumerate}
\end{definition}

\begin{example}
\label{ex:vector_spaces}
Standard examples:
\begin{enumerate}[nosep]
    \item $\F^n = \{(x_1, \ldots, x_n) : x_i \in \F\}$ with componentwise operations.
    \item $\F[x]_{\le n}$, the polynomials of degree at most $n$, with the usual operations.
    \item $C[a,b]$, the continuous functions $f: [a,b] \to \R$, with pointwise operations.
    \item $\R^{m \times n}$, the $m \times n$ matrices with real entries.
\end{enumerate}
\end{example}

\subsection{Subspaces}
\label{subsec:subspaces}

\begin{definition}[Subspace]
\label{def:subspace}
A subset $U \subseteq V$ is a \emph{subspace} if $U$ is itself a vector space under
the same operations. Equivalently, $U$ is a subspace if and only if:
\begin{enumerate}[nosep]
    \item $0 \in U$,
    \item $u + v \in U$ for all $u, v \in U$ (closed under addition),
    \item $\lambda u \in U$ for all $\lambda \in \F$, $u \in U$ (closed under scalar multiplication).
\end{enumerate}
\end{definition}

\begin{proposition}[Intersection of Subspaces]
\label{prop:intersection_subspaces}
If $U_1, U_2$ are subspaces of $V$, then $U_1 \cap U_2$ is a subspace of $V$.
\end{proposition}

\begin{proof}
We verify the three conditions. (1) $0 \in U_1$ and $0 \in U_2$, so $0 \in U_1 \cap U_2$.
(2) If $u, v \in U_1 \cap U_2$, then $u, v \in U_1$ so $u + v \in U_1$, and similarly
$u + v \in U_2$, so $u + v \in U_1 \cap U_2$. (3) If $u \in U_1 \cap U_2$ and $\lambda \in \F$,
then $\lambda u \in U_1$ and $\lambda u \in U_2$, so $\lambda u \in U_1 \cap U_2$.
\end{proof}

\subsection{Span and Linear Independence}
\label{subsec:span_independence}

\begin{definition}[Span]
\label{def:span}
The \emph{span} of a list of vectors $v_1, \ldots, v_m \in V$ is
\[
    \Span(v_1, \ldots, v_m) = \{\lambda_1 v_1 + \cdots + \lambda_m v_m : \lambda_1, \ldots, \lambda_m \in \F\}.
\]
\end{definition}

\begin{definition}[Linear Independence]
\label{def:linear_independence}
Vectors $v_1, \ldots, v_m \in V$ are \emph{linearly independent} if the only solution to
$\lambda_1 v_1 + \cdots + \lambda_m v_m = 0$ is $\lambda_1 = \cdots = \lambda_m = 0$.
They are \emph{linearly dependent} otherwise.
\end{definition}

\begin{lemma}[Linear Dependence Lemma]
\label{lem:lin_dep}
If $v_1, \ldots, v_m$ are linearly dependent and $v_1 \neq 0$, then there exists
$j \in \{2, \ldots, m\}$ such that $v_j \in \Span(v_1, \ldots, v_{j-1})$ and
\[
    \Span(v_1, \ldots, v_{j-1}, v_{j+1}, \ldots, v_m) = \Span(v_1, \ldots, v_m).
\]
\end{lemma}

\begin{proof}
Since the vectors are linearly dependent, there exist $\lambda_1, \ldots, \lambda_m$,
not all zero, with $\sum_{i=1}^m \lambda_i v_i = 0$. Let $j$ be the largest index with
$\lambda_j \neq 0$. Then $j \geq 2$ (since $v_1 \neq 0$ means $\lambda_1 = 0$ would
force all $\lambda_i = 0$ if $j = 1$). We have
$v_j = -\lambda_j^{-1} \sum_{i=1}^{j-1} \lambda_i v_i \in \Span(v_1, \ldots, v_{j-1})$.
Any vector in $\Span(v_1, \ldots, v_m)$ can replace its $v_j$-component using this
expression, so removing $v_j$ does not change the span.
\end{proof}

\subsection{Basis and Dimension}
\label{subsec:basis_dimension}

\begin{definition}[Basis]
\label{def:basis}
A \emph{basis} of $V$ is a list of vectors in $V$ that is linearly independent and spans $V$.
\end{definition}

\begin{theorem}[Bases Have the Same Length]
\label{thm:dimension}
If $V$ is finite-dimensional, then any two bases of $V$ have the same number of elements.
\end{theorem}

\begin{proof}
Let $\{u_1, \ldots, u_m\}$ and $\{w_1, \ldots, w_n\}$ be two bases. We show $m \leq n$.
Since the $w$'s span $V$, the list $w_1, \ldots, w_n, u_1$ is linearly dependent. By the
linear dependence lemma, we can remove some $w_j$ while preserving the span. Repeating this
process, each time we add a $u_i$ and remove a $w_j$, after $m$ steps we have a spanning list
consisting of all $m$ vectors $u_i$ and $n - m$ remaining $w$'s. This requires $n - m \geq 0$,
so $m \leq n$. By symmetry, $n \leq m$, hence $m = n$.
\end{proof}

\begin{definition}[Dimension]
\label{def:dimension}
The \emph{dimension} of a finite-dimensional vector space $V$, denoted $\dim V$, is the
number of elements in any basis of $V$.
\end{definition}

\begin{proposition}
\label{prop:basis_extension}
Every linearly independent list in a finite-dimensional vector space can be extended to a basis.
\end{proposition}

\begin{proof}
Let $v_1, \ldots, v_m$ be linearly independent in $V$ with $\dim V = n$. If $m = n$, we
are done (a linearly independent list of length $n$ in an $n$-dimensional space must span).
If $m < n$, the list does not span $V$, so there exists $v_{m+1} \notin \Span(v_1, \ldots, v_m)$.
Then $v_1, \ldots, v_{m+1}$ is linearly independent. Repeat until the list has length $n$.
\end{proof}

%% ============================================================
%% SECTION 3: LINEAR MAPS
%% ============================================================
\section{Linear Maps}
\label{sec:linear_maps}

\begin{definition}[Linear Map]
\label{def:linear_map}
A function $T: V \to W$ between vector spaces over $\F$ is a \emph{linear map} if
$T(u + v) = Tu + Tv$ and $T(\lambda v) = \lambda Tv$ for all $u, v \in V$ and $\lambda \in \F$.
\end{definition}

\begin{definition}[Null Space and Range]
\label{def:null_range}
For a linear map $T: V \to W$:
\begin{itemize}[nosep]
    \item The \emph{null space} (kernel) is $\nullspace T = \{v \in V : Tv = 0\}$.
    \item The \emph{range} (image) is $\range T = \{Tv : v \in V\}$.
\end{itemize}
\end{definition}

\begin{proposition}
\label{prop:injectivity}
A linear map $T$ is injective if and only if $\nullspace T = \{0\}$.
\end{proposition}

\begin{proof}
($\Rightarrow$) If $T$ is injective and $Tv = 0 = T0$, then $v = 0$.
($\Leftarrow$) If $\nullspace T = \{0\}$ and $Tu = Tv$, then $T(u - v) = 0$,
so $u - v = 0$, hence $u = v$.
\end{proof}

\begin{theorem}[Rank-Nullity Theorem]
\label{thm:rank_nullity}
If $V$ is finite-dimensional and $T: V \to W$ is linear, then
\[
    \dim V = \dim(\nullspace T) + \dim(\range T).
\]
\end{theorem}

\begin{proof}
Let $u_1, \ldots, u_m$ be a basis of $\nullspace T$ and extend it to a basis
$u_1, \ldots, u_m, v_1, \ldots, v_k$ of $V$. We claim $Tv_1, \ldots, Tv_k$ is a basis of
$\range T$.

\textbf{Spanning:} Any $w \in \range T$ has $w = T(\lambda_1 u_1 + \cdots + \lambda_m u_m
+ \mu_1 v_1 + \cdots + \mu_k v_k) = \mu_1 Tv_1 + \cdots + \mu_k Tv_k$ (since $Tu_i = 0$).

\textbf{Independence:} If $\mu_1 Tv_1 + \cdots + \mu_k Tv_k = 0$, then
$T(\mu_1 v_1 + \cdots + \mu_k v_k) = 0$, so $\mu_1 v_1 + \cdots + \mu_k v_k \in \nullspace T$.
Thus $\mu_1 v_1 + \cdots + \mu_k v_k = \alpha_1 u_1 + \cdots + \alpha_m u_m$ for some scalars.
Since the full list is a basis, all coefficients are zero, so $\mu_i = 0$ for all $i$.

Therefore $\dim(\range T) = k$ and $\dim V = m + k = \dim(\nullspace T) + \dim(\range T)$.
\end{proof}

\subsection{Matrix Representation}
\label{subsec:matrix_rep}

Fix bases $\{v_1, \ldots, v_n\}$ of $V$ and $\{w_1, \ldots, w_m\}$ of $W$. Any linear map
$T: V \to W$ is uniquely determined by its values on the basis:
$Tv_j = \sum_{i=1}^m A_{ij} w_i$. The $m \times n$ matrix $A = (A_{ij})$ is the
\emph{matrix of $T$} with respect to the chosen bases.

\begin{proposition}[Isomorphism of Finite-Dimensional Spaces]
\label{prop:isomorphism}
Two finite-dimensional vector spaces over $\F$ are isomorphic if and only if they have the
same dimension.
\end{proposition}

\begin{proof}
($\Rightarrow$) If $T: V \to W$ is an isomorphism and $\{v_1, \ldots, v_n\}$ is a basis
of $V$, then $\{Tv_1, \ldots, Tv_n\}$ is a basis of $W$ (injectivity gives independence,
surjectivity gives spanning), so $\dim W = n = \dim V$.

($\Leftarrow$) If $\dim V = \dim W = n$, pick bases $\{v_i\}$ of $V$ and $\{w_i\}$ of $W$.
Define $T: V \to W$ by $Tv_i = w_i$ and extend linearly. This is well-defined, linear,
and bijective (it maps a basis to a basis).
\end{proof}

\begin{rlconnection}[Bellman Equation as a Linear System]
For a fixed policy $\pi$ in a finite MDP, the Bellman equation $V^\pi = r^\pi + \gamma P^\pi V^\pi$
is a linear system $(I - \gamma P^\pi) V^\pi = r^\pi$. It has a unique solution because
$I - \gamma P^\pi$ is invertible (its eigenvalues are $1 - \gamma \lambda_i$ where $|\lambda_i| \leq 1$
and $\gamma < 1$, so all eigenvalues are nonzero).
\end{rlconnection}

%% ============================================================
%% SECTION 4: DETERMINANTS
%% ============================================================
\section{Determinants}
\label{sec:determinants}

We define the determinant for $n \times n$ matrices over $\F$ and establish the key properties.

\begin{definition}[Determinant]
\label{def:determinant}
For $A \in \F^{n \times n}$, the \emph{determinant} is defined recursively by cofactor
expansion along the first row:
\[
    \det A = \sum_{j=1}^n (-1)^{1+j} A_{1j} \det A_{1j}',
\]
where $A_{1j}'$ is the $(n-1) \times (n-1)$ matrix obtained by deleting row~1 and column~$j$.
The base case is $\det(a) = a$ for $1 \times 1$ matrices.
\end{definition}

\begin{theorem}[Properties of the Determinant]
\label{thm:det_properties}
For $A, B \in \F^{n \times n}$:
\begin{enumerate}[nosep]
    \item $\det I = 1$.
    \item The determinant is multilinear and alternating in the rows of $A$.
    \item If $A$ has two equal rows, then $\det A = 0$.
    \item $\det A^T = \det A$.
    \item Row operations: swapping two rows multiplies $\det$ by $-1$; adding a scalar multiple
          of one row to another leaves $\det$ unchanged.
\end{enumerate}
\end{theorem}

\begin{proof}
These follow from the recursive definition and induction on $n$. We prove (3) and (5).

\textbf{(3)} If rows $p$ and $q$ (with $p \neq q$) are equal, swapping them leaves $A$
unchanged. By the alternating property (from multilinearity applied to the sum of two
equal rows), $\det A = -\det A$, so $\det A = 0$.

\textbf{(5)} Adding $c$ times row $p$ to row $q$ ($p \neq q$): by multilinearity in row $q$,
$\det A' = \det A + c \cdot \det B$ where $B$ is the matrix with row $q$ replaced by row $p$.
Since $B$ has two equal rows, $\det B = 0$ by (3).
\end{proof}

\begin{theorem}[Multiplicativity]
\label{thm:det_multiplicativity}
For $A, B \in \F^{n \times n}$, $\det(AB) = \det A \cdot \det B$.
\end{theorem}

\begin{proof}
If $A$ is not invertible, then $\rank(AB) \leq \rank(A) < n$, so $AB$ is not invertible and
both sides are zero. If $A$ is invertible, it is a product of elementary matrices (row operations).
For elementary matrices $E$, one verifies directly that $\det(EB) = \det(E) \det(B)$: a row
swap gives $\det(E) = -1$, a row scaling by $c$ gives $\det(E) = c$, and a row addition gives
$\det(E) = 1$. The result follows by induction on the number of elementary factors.
\end{proof}

\begin{corollary}
\label{cor:det_invertibility}
$A$ is invertible if and only if $\det A \neq 0$.
\end{corollary}

\begin{proof}
($\Rightarrow$) If $A$ is invertible, $1 = \det I = \det(A A^{-1}) = \det A \cdot \det A^{-1}$,
so $\det A \neq 0$.
($\Leftarrow$) If $A$ is not invertible, row reduction produces a row of zeros, giving $\det A = 0$.
\end{proof}

%% ============================================================
%% SECTION 5: EIGENVALUES AND EIGENVECTORS
%% ============================================================
\section{Eigenvalues and Eigenvectors}
\label{sec:eigenvalues}

\begin{definition}[Eigenvalue, Eigenvector]
\label{def:eigenvalue}
Let $T: V \to V$ be a linear operator on a finite-dimensional space $V$. A scalar
$\lambda \in \F$ is an \emph{eigenvalue} of $T$ if there exists a nonzero $v \in V$ with
$Tv = \lambda v$. Such a $v$ is an \emph{eigenvector} corresponding to $\lambda$.
\end{definition}

\begin{definition}[Characteristic Polynomial]
\label{def:char_poly}
For $A \in \F^{n \times n}$, the \emph{characteristic polynomial} is
$p(\lambda) = \det(\lambda I - A)$.
\end{definition}

\begin{theorem}
\label{thm:eigenvalues_are_roots}
$\lambda$ is an eigenvalue of $A$ if and only if $\det(\lambda I - A) = 0$.
\end{theorem}

\begin{proof}
$\lambda$ is an eigenvalue $\iff$ there exists $v \neq 0$ with $Av = \lambda v$ $\iff$
$(\lambda I - A)v = 0$ has a nontrivial solution $\iff$ $\lambda I - A$ is not invertible
$\iff$ $\det(\lambda I - A) = 0$.
\end{proof}

\begin{theorem}[Eigenvectors for Distinct Eigenvalues]
\label{thm:distinct_eigenvalues}
If $\lambda_1, \ldots, \lambda_m$ are distinct eigenvalues of $T$ with corresponding
eigenvectors $v_1, \ldots, v_m$, then $v_1, \ldots, v_m$ are linearly independent.
\end{theorem}

\begin{proof}
By induction on $m$. The case $m = 1$ is clear ($v_1 \neq 0$). Suppose the result holds for
$m - 1$ eigenvectors and assume $\sum_{i=1}^m c_i v_i = 0$. Apply $T$ to get
$\sum_{i=1}^m c_i \lambda_i v_i = 0$. Subtract $\lambda_m$ times the original equation:
$\sum_{i=1}^{m-1} c_i (\lambda_i - \lambda_m) v_i = 0$. By the inductive hypothesis,
$v_1, \ldots, v_{m-1}$ are linearly independent, so $c_i(\lambda_i - \lambda_m) = 0$ for
$i = 1, \ldots, m-1$. Since the eigenvalues are distinct, $\lambda_i - \lambda_m \neq 0$,
so $c_i = 0$ for $i < m$. The original equation then gives $c_m v_m = 0$, so $c_m = 0$.
\end{proof}

\begin{corollary}[Sufficient Condition for Diagonalizability]
\label{cor:diag_distinct}
If $A \in \F^{n \times n}$ has $n$ distinct eigenvalues, then $A$ is diagonalizable.
\end{corollary}

\begin{proof}
By Theorem~\ref{thm:distinct_eigenvalues}, the $n$ eigenvectors are linearly independent and
thus form a basis. If $P$ is the matrix whose columns are these eigenvectors, then
$P^{-1}AP = \diag(\lambda_1, \ldots, \lambda_n)$.
\end{proof}

\begin{rlconnection}[Eigenvalues and Convergence Rates]
Eigenvalue analysis is central to understanding convergence rates. The spectral radius of
a transition matrix $P$ determines the mixing time of a Markov chain. The condition number
$\kappa = \lambda_{\max}/\lambda_{\min}$ of the Hessian determines gradient descent convergence
speed: gradient descent on a quadratic with condition number $\kappa$ requires
$O(\kappa \log(1/\varepsilon))$ iterations.
\end{rlconnection}

%% ============================================================
%% SECTION 6: INNER PRODUCT SPACES
%% ============================================================
\section{Inner Product Spaces}
\label{sec:inner_product}

\begin{definition}[Inner Product]
\label{def:inner_product}
An \emph{inner product} on a real vector space $V$ is a function
$\langle \cdot, \cdot \rangle: V \times V \to \R$ satisfying:
\begin{enumerate}[nosep]
    \item \textbf{Symmetry:} $\langle u, v \rangle = \langle v, u \rangle$.
    \item \textbf{Linearity in the first argument:} $\langle \alpha u + \beta v, w \rangle
          = \alpha \langle u, w \rangle + \beta \langle v, w \rangle$.
    \item \textbf{Positive definiteness:} $\langle v, v \rangle \geq 0$ with equality iff $v = 0$.
\end{enumerate}
The pair $(V, \langle \cdot, \cdot \rangle)$ is an \emph{inner product space}.
\end{definition}

\begin{remark}
The standard inner product on $\R^n$ is $\langle x, y \rangle = \sum_{i=1}^n x_i y_i = x^T y$.
\end{remark}

\begin{theorem}[Cauchy--Schwarz Inequality]
\label{thm:cauchy_schwarz}
For all $u, v$ in an inner product space,
\[
    |\langle u, v \rangle| \leq \|u\| \cdot \|v\|,
\]
where $\|w\| = \sqrt{\langle w, w \rangle}$. Equality holds if and only if $u$ and $v$
are linearly dependent.
\end{theorem}

\begin{proof}
If $v = 0$, both sides are zero. Assume $v \neq 0$ and let $t \in \R$. Then
\[
    0 \leq \|u - tv\|^2 = \langle u - tv, u - tv \rangle = \|u\|^2 - 2t\langle u, v\rangle + t^2\|v\|^2.
\]
This is a nonnegative quadratic in $t$, so its discriminant is non-positive:
$4\langle u, v \rangle^2 - 4\|u\|^2\|v\|^2 \leq 0$, giving
$|\langle u, v \rangle| \leq \|u\| \|v\|$.

Equality holds iff the discriminant is zero, i.e., the minimum of the quadratic is zero,
meaning $u - tv = 0$ for some $t$, i.e., $u$ and $v$ are linearly dependent.
\end{proof}

\subsection{Orthogonality and Gram--Schmidt}
\label{subsec:gram_schmidt}

\begin{definition}[Orthogonal, Orthonormal]
\label{def:orthogonal}
Vectors $u, v$ are \emph{orthogonal} if $\langle u, v \rangle = 0$. A list
$e_1, \ldots, e_m$ is \emph{orthonormal} if $\langle e_i, e_j \rangle = \delta_{ij}$
(the Kronecker delta).
\end{definition}

\begin{proposition}
\label{prop:orthonormal_independent}
An orthonormal list is linearly independent.
\end{proposition}

\begin{proof}
If $\sum c_i e_i = 0$, take the inner product with $e_j$:
$0 = \langle \sum c_i e_i, e_j \rangle = \sum c_i \langle e_i, e_j \rangle = c_j$.
\end{proof}

\begin{theorem}[Gram--Schmidt Process]
\label{thm:gram_schmidt}
Given a linearly independent list $v_1, \ldots, v_m$ in an inner product space, there exists
an orthonormal list $e_1, \ldots, e_m$ such that
$\Span(e_1, \ldots, e_k) = \Span(v_1, \ldots, v_k)$ for each $k = 1, \ldots, m$.
\end{theorem}

\begin{proof}
\textbf{Algorithm:} Set $e_1 = v_1 / \|v_1\|$. For $k = 2, \ldots, m$, define
\[
    w_k = v_k - \sum_{j=1}^{k-1} \langle v_k, e_j \rangle e_j, \qquad e_k = \frac{w_k}{\|w_k\|}.
\]

\textbf{Well-defined:} $w_k \neq 0$ because $v_k \notin \Span(v_1, \ldots, v_{k-1})
= \Span(e_1, \ldots, e_{k-1})$ by linear independence.

\textbf{Orthogonality:} For $j < k$,
$\langle w_k, e_j \rangle = \langle v_k, e_j \rangle - \langle v_k, e_j \rangle = 0$.

\textbf{Span preservation:} By construction, $e_k \in \Span(v_1, \ldots, v_k)$ and
$v_k = w_k + \sum \langle v_k, e_j \rangle e_j \in \Span(e_1, \ldots, e_k)$, so the
spans agree at each step.
\end{proof}

\subsection{Orthogonal Projection}
\label{subsec:projection}

\begin{theorem}[Best Approximation]
\label{thm:best_approx}
Let $U$ be a finite-dimensional subspace of an inner product space $V$ and let $v \in V$.
There exists a unique $u_0 \in U$ minimizing $\|v - u\|$ over $u \in U$, given by the
\emph{orthogonal projection}:
\[
    u_0 = \text{proj}_U(v) = \sum_{i=1}^m \langle v, e_i \rangle e_i,
\]
where $e_1, \ldots, e_m$ is any orthonormal basis of $U$. Moreover, $v - u_0 \perp U$.
\end{theorem}

\begin{proof}
Define $u_0 = \sum \langle v, e_i \rangle e_i$. For any $e_j$,
$\langle v - u_0, e_j \rangle = \langle v, e_j \rangle - \langle v, e_j \rangle = 0$,
so $v - u_0 \perp U$.

For any $u \in U$, write $u - u_0 \in U$ and note $(v - u_0) \perp (u_0 - u)$:
\[
    \|v - u\|^2 = \|(v - u_0) + (u_0 - u)\|^2 = \|v - u_0\|^2 + \|u_0 - u\|^2 \geq \|v - u_0\|^2,
\]
with equality iff $u = u_0$.
\end{proof}

\begin{keyresult}[Orthogonal Projection: The Best Approximation]
The orthogonal projection $\text{proj}_U(v)$ is the unique element of $U$ closest to $v$.
This is the geometric essence of least squares: given data and a model subspace, the least
squares solution is the orthogonal projection of the data onto the model space. This principle
generalizes to Hilbert spaces in functional analysis (Phase~01).
\end{keyresult}

%% ============================================================
%% SECTION 7: THE SPECTRAL THEOREM
%% ============================================================
\section{The Spectral Theorem}
\label{sec:spectral}

We now restrict to real symmetric matrices, which are the most important case for
optimization and machine learning.

\begin{theorem}[Real Eigenvalues of Symmetric Matrices]
\label{thm:symmetric_real_eigenvalues}
Every real symmetric matrix has only real eigenvalues.
\end{theorem}

\begin{proof}
Let $A \in \R^{n \times n}$ be symmetric with $Av = \lambda v$ for some $v \in \C^n$,
$v \neq 0$. Consider $\bar{v}^T A v$. On one hand,
$\bar{v}^T A v = \bar{v}^T (\lambda v) = \lambda \bar{v}^T v = \lambda \sum |v_i|^2$.
On the other hand, since $A = A^T$ is real,
$\bar{v}^T A v = (A^T \bar{v})^T v = (A\bar{v})^T v = (\bar{\lambda}\bar{v})^T v
= \bar{\lambda} \bar{v}^T v = \bar{\lambda} \sum |v_i|^2$.
Since $\sum |v_i|^2 > 0$, we get $\lambda = \bar{\lambda}$, so $\lambda \in \R$.
\end{proof}

\begin{theorem}[Orthogonality of Eigenvectors]
\label{thm:symmetric_orthogonal_eigenvectors}
Eigenvectors of a real symmetric matrix corresponding to distinct eigenvalues are orthogonal.
\end{theorem}

\begin{proof}
Let $Av = \lambda v$ and $Aw = \mu w$ with $\lambda \neq \mu$. Then
$\lambda \langle v, w \rangle = \langle Av, w \rangle = \langle v, A^T w \rangle
= \langle v, Aw \rangle = \mu \langle v, w \rangle$.
So $(\lambda - \mu)\langle v, w \rangle = 0$. Since $\lambda \neq \mu$,
$\langle v, w \rangle = 0$.
\end{proof}

\begin{theorem}[Real Spectral Theorem]
\label{thm:spectral}
Every real symmetric matrix $A \in \R^{n \times n}$ is orthogonally diagonalizable:
there exists an orthogonal matrix $Q$ (i.e., $Q^T Q = I$) and a diagonal matrix
$\Lambda = \diag(\lambda_1, \ldots, \lambda_n)$ such that $A = Q\Lambda Q^T$.
\end{theorem}

\begin{proof}
We proceed by induction on $n$. The case $n = 1$ is trivial.

Assume the result holds for $(n-1) \times (n-1)$ symmetric matrices. By
Theorem~\ref{thm:symmetric_real_eigenvalues}, $A$ has a real eigenvalue $\lambda_1$ with
a unit eigenvector $v_1 \in \R^n$. Extend $v_1$ to an orthonormal basis
$\{v_1, v_2, \ldots, v_n\}$ of $\R^n$ (using Gram--Schmidt on any extension).
Let $P = [v_1 \mid v_2 \mid \cdots \mid v_n]$. Then
\[
    P^T A P = \begin{pmatrix} \lambda_1 & b^T \\ b & B \end{pmatrix},
\]
where $b \in \R^{n-1}$ and $B \in \R^{(n-1) \times (n-1)}$. Since $P^TAP$ is symmetric
(as $A$ is), the block structure gives $b^T = b^T$ (automatically) and we need $b = 0$.
Indeed, the first column of $P^TAP$ is $P^TAv_1 = P^T(\lambda_1 v_1) = \lambda_1 e_1$,
so $b = 0$. Also $B$ is symmetric.

By the inductive hypothesis, $B = Q' \Lambda' (Q')^T$ for some orthogonal $Q'$ and diagonal
$\Lambda'$. Set $Q_0 = \begin{pmatrix} 1 & 0 \\ 0 & Q' \end{pmatrix}$ and
$Q = PQ_0$. Then $Q$ is orthogonal and
$Q^T A Q = Q_0^T (P^T A P) Q_0 = \diag(\lambda_1, \Lambda') = \Lambda$.
\end{proof}

\begin{definition}[Positive Definite and Positive Semidefinite]
\label{def:positive_definite}
A symmetric matrix $A$ is:
\begin{itemize}[nosep]
    \item \emph{Positive definite} ($A \succ 0$) if $x^T A x > 0$ for all $x \neq 0$.
    \item \emph{Positive semidefinite} ($A \succeq 0$) if $x^T A x \geq 0$ for all $x$.
\end{itemize}
\end{definition}

\begin{theorem}[Eigenvalue Characterization]
\label{thm:pd_eigenvalues}
A symmetric matrix $A$ is positive definite if and only if all its eigenvalues are positive.
It is positive semidefinite if and only if all its eigenvalues are nonnegative.
\end{theorem}

\begin{proof}
By the spectral theorem, $A = Q\Lambda Q^T$. For any $x \neq 0$, let $y = Q^T x \neq 0$.
Then $x^T A x = y^T \Lambda y = \sum \lambda_i y_i^2$. This is positive for all $y \neq 0$
iff all $\lambda_i > 0$, and nonnegative for all $y$ iff all $\lambda_i \geq 0$.
\end{proof}

\begin{keyresult}[The Spectral Theorem in Practice]
Every real symmetric matrix $A$ can be written as $A = Q\Lambda Q^T$ where $Q$ is orthogonal
and $\Lambda$ is diagonal. This decomposition is the foundation of:
\begin{enumerate}[nosep]
    \item \textbf{PCA} --- the principal components are eigenvectors of the covariance matrix.
    \item \textbf{Convexity} --- a twice-differentiable function is convex iff its Hessian is
          positive semidefinite (all eigenvalues $\geq 0$).
    \item \textbf{Natural policy gradient} --- the update direction is $F^{-1}\nabla J$,
          where $F$ is the (symmetric, positive semidefinite) Fisher information matrix.
\end{enumerate}
\end{keyresult}

%% ============================================================
%% SECTION 8: SINGULAR VALUE DECOMPOSITION
%% ============================================================
\section{Singular Value Decomposition}
\label{sec:svd}

The SVD extends diagonalization to rectangular matrices.

\begin{theorem}[Singular Value Decomposition]
\label{thm:svd}
Every matrix $A \in \R^{m \times n}$ can be written as
\[
    A = U \Sigma V^T,
\]
where $U \in \R^{m \times m}$ is orthogonal, $V \in \R^{n \times n}$ is orthogonal, and
$\Sigma \in \R^{m \times n}$ is ``diagonal'' with nonnegative entries
$\sigma_1 \geq \sigma_2 \geq \cdots \geq \sigma_{\min(m,n)} \geq 0$ (the \emph{singular values}).
\end{theorem}

\begin{proof}
The matrix $A^TA \in \R^{n \times n}$ is symmetric and positive semidefinite (since
$x^T A^T A x = \|Ax\|^2 \geq 0$). By the spectral theorem, $A^T A = V D V^T$ where
$D = \diag(\sigma_1^2, \ldots, \sigma_n^2)$ with $\sigma_i \geq 0$ and $V$ orthogonal.

For each $\sigma_i > 0$, define $u_i = Av_i / \sigma_i$. Then the $u_i$ are orthonormal:
$\langle u_i, u_j \rangle = \frac{1}{\sigma_i \sigma_j} v_i^T A^T A v_j
= \frac{\sigma_j^2}{\sigma_i \sigma_j} v_i^T v_j = \delta_{ij}$.

Extend $\{u_1, \ldots, u_r\}$ (where $r = \rank A$) to an orthonormal basis of $\R^m$.
Forming $U = [u_1 \mid \cdots \mid u_m]$ gives $AV = U\Sigma$, hence $A = U\Sigma V^T$.
\end{proof}

\begin{remark}[Geometric Interpretation]
The SVD says: any linear map $\R^n \to \R^m$ can be decomposed as (1) rotate in the domain
(apply $V^T$), (2) scale along coordinate axes (apply $\Sigma$), (3) rotate in the codomain
(apply $U$). The singular values are the semi-axis lengths of the image of the unit ball.
\end{remark}

\begin{theorem}[Eckart--Young: Best Low-Rank Approximation]
\label{thm:eckart_young}
Let $A = U\Sigma V^T$ be the SVD of $A$ with singular values $\sigma_1 \geq \cdots \geq \sigma_r > 0$.
For any $k < r$, the matrix $A_k = \sum_{i=1}^k \sigma_i u_i v_i^T$ minimizes $\|A - B\|_F$
over all matrices $B$ of rank at most $k$, and $\|A - A_k\|_F = \sqrt{\sigma_{k+1}^2 + \cdots + \sigma_r^2}$.
\end{theorem}

\begin{proof}
Write $A = \sum_{i=1}^r \sigma_i u_i v_i^T$. Since the $u_i v_i^T$ have orthonormal columns
and rows, $\|A\|_F^2 = \sum \sigma_i^2$ (the Frobenius norm depends only on singular values).
For $A_k = \sum_{i=1}^k \sigma_i u_i v_i^T$,
$\|A - A_k\|_F^2 = \sum_{i=k+1}^r \sigma_i^2$.

To show this is minimal, let $B$ have rank $\leq k$, so $\nullspace B$ has dimension $\geq n - k$.
The space $\Span(v_1, \ldots, v_{k+1})$ has dimension $k+1$, so it intersects
$\nullspace B$ nontrivially. Pick a unit $z$ in the intersection. Then
$\|A - B\|_F^2 \geq \|(A-B)z\|^2 = \|Az\|^2 \geq \sigma_{k+1}^2$
(since $z$ lies in the span of the top $k+1$ right singular vectors). Repeating this
argument with an appropriate subspace gives $\|A-B\|_F^2 \geq \sum_{i=k+1}^r \sigma_i^2$.
\end{proof}

\begin{rlconnection}[SVD in Machine Learning]
SVD is the workhorse of dimensionality reduction and matrix approximation. In RL,
low-rank approximations appear in linear function approximation (projecting value functions
onto a low-dimensional feature space) and in compressing large state-action value tables.
PCA is a direct application: the principal components are the right singular vectors of the
centered data matrix.
\end{rlconnection}

%% ============================================================
%% SECTION 9: MATRIX NORMS AND THE MATRIX EXPONENTIAL
%% ============================================================
\section{Matrix Norms and the Matrix Exponential}
\label{sec:matrix_norms}

\begin{definition}[Operator Norm]
\label{def:operator_norm}
For $A \in \R^{m \times n}$ and a vector norm $\|\cdot\|$ on both $\R^n$ and $\R^m$,
the \emph{operator norm} (induced norm) is
\[
    \|A\|_{\text{op}} = \sup_{x \neq 0} \frac{\|Ax\|}{\|x\|} = \sup_{\|x\| = 1} \|Ax\|.
\]
For the Euclidean norm, $\|A\|_{\text{op}} = \sigma_1(A)$ (the largest singular value).
\end{definition}

\begin{definition}[Frobenius Norm]
\label{def:frobenius}
$\|A\|_F = \sqrt{\sum_{i,j} A_{ij}^2} = \sqrt{\tr(A^T A)} = \sqrt{\sum_i \sigma_i^2}$.
\end{definition}

\begin{theorem}[Submultiplicativity]
\label{thm:submultiplicativity}
For the operator norm, $\|AB\|_{\text{op}} \leq \|A\|_{\text{op}} \|B\|_{\text{op}}$.
\end{theorem}

\begin{proof}
For any $x$ with $\|x\| = 1$, $\|ABx\| \leq \|A\|_{\text{op}} \|Bx\| \leq \|A\|_{\text{op}} \|B\|_{\text{op}} \|x\| = \|A\|_{\text{op}} \|B\|_{\text{op}}$. Taking the supremum over $x$ gives the result.
\end{proof}

\subsection{The Matrix Exponential}
\label{subsec:matrix_exp}

\begin{definition}[Matrix Exponential]
\label{def:matrix_exp}
For $A \in \R^{n \times n}$, the \emph{matrix exponential} is
\[
    e^A = \sum_{k=0}^{\infty} \frac{A^k}{k!},
\]
which converges absolutely in any matrix norm (since $\sum \|A\|^k / k! = e^{\|A\|} < \infty$).
\end{definition}

\begin{proposition}[Properties of the Matrix Exponential]
\label{prop:matrix_exp_properties}
\begin{enumerate}[nosep]
    \item $e^{0} = I$.
    \item If $AB = BA$, then $e^{A+B} = e^A e^B$.
    \item $e^A$ is always invertible, with $(e^A)^{-1} = e^{-A}$.
    \item $\frac{d}{dt} e^{tA} = A e^{tA}$.
\end{enumerate}
\end{proposition}

\begin{proof}
\textbf{(1)} $e^0 = \sum 0^k/k! = I$ (the $k=0$ term).

\textbf{(2)} When $AB = BA$, the binomial theorem gives
$(A+B)^k = \sum_{j=0}^k \binom{k}{j} A^j B^{k-j}$. Then
$e^{A+B} = \sum_k \frac{1}{k!} \sum_j \binom{k}{j} A^j B^{k-j}
= \sum_j \sum_\ell \frac{A^j}{j!} \frac{B^\ell}{\ell!} = e^A e^B$
(reindexing $\ell = k - j$ and using absolute convergence to rearrange).

\textbf{(3)} Since $A$ and $-A$ commute, $e^A e^{-A} = e^{A-A} = e^0 = I$.

\textbf{(4)} $\frac{d}{dt} e^{tA} = \frac{d}{dt} \sum_k \frac{(tA)^k}{k!}
= \sum_{k=1}^\infty \frac{k t^{k-1} A^k}{k!} = A \sum_{k=1}^\infty \frac{(tA)^{k-1}}{(k-1)!}
= A e^{tA}$.
\end{proof}

\begin{proposition}[Matrix Exponential of Diagonalizable Matrices]
\label{prop:exp_diag}
If $A = P D P^{-1}$ with $D = \diag(\lambda_1, \ldots, \lambda_n)$, then
$e^A = P \diag(e^{\lambda_1}, \ldots, e^{\lambda_n}) P^{-1}$.
\end{proposition}

\begin{proof}
$A^k = P D^k P^{-1}$, so $e^A = \sum_k P D^k P^{-1} / k! = P (\sum_k D^k / k!) P^{-1}
= P e^D P^{-1}$. Since $D$ is diagonal, $e^D = \diag(e^{\lambda_1}, \ldots, e^{\lambda_n})$.
\end{proof}

\begin{rlconnection}[Matrix Exponential in Dynamical Systems]
The matrix exponential solves the linear ODE $\dot{x}(t) = Ax(t)$: the solution is
$x(t) = e^{tA} x(0)$. This appears in continuous-time Markov chains (where the transition
semigroup is $P(t) = e^{tQ}$ for a rate matrix $Q$) and in the analysis of linear dynamical
systems that arise in control theory (Phase~06).
\end{rlconnection}

\end{document}
